% !TeX root = ../main.tex

\section{Empirical Study}\label{motivation}
 

\subsection{Data Collection}
We collect our dataset through a two-step process. First, we identify CWE entries associated with overflow vulnerabilities. Then, we extract relevant vulnerability samples from widely-used vulnerability repair datasets based on these identified CWEs.

To systematically identify overflow-related vulnerabilities, we adopt a two-pronged approach to collecting relevant CWE entries. First, we refer to the Common Weakness Enumeration (CWE) framework\cite{CWEHome}, specifically utilizing CWE View-1000\cite{CWE1000}, also known as the “Research Concepts View.” This view provides a semantically organized structure of software weaknesses tailored for academic analysis, emphasizing conceptual clarity and hierarchy. Based on this taxonomy, we identify CWE-119 — Improper Restriction of Operations within the Bounds of a Memory Buffer — and its descendant nodes as representative of overflow vulnerabilities. The complete list of selected CWE entries is provided in Table~\ref{tab:cwe-list}.Second, to complement the taxonomy-driven selection, we conduct a survey of recent research on vulnerability detection and repair. Using keywords such as “overflow,” “detect,” and “repair,” we search for papers published in top-tier software engineering and security venues (e.g., CCS\footnote{ACM Conference on Computer and Communications Security}, ASE\footnote{IEEE/ACM International Conference on Automated Software Engineering}). From these works\cite{le2024study,li2022pacmem,liu2025cctag,zhang2023vulnerability,kang2022tracer,gorter2024sticky,steenhoek2023empirical,ba2022efficient,guo2024precise}, we analyze the datasets used and extract the CWE identifiers associated with overflow vulnerabilities. We find that most of these CWEs are already included in the CWE View-1000 hierarchy; for those not covered, we incorporate them into our list to ensure comprehensive coverage and consistency with real-world research practices.

Next, we extract overflow-related vulnerability samples from two widely used vulnerability repair datasets: BigVul\cite{bigvul} and CVEFixes\cite{bhandari2021:cvefixes}. BigVul is a large-scale dataset that collects real-world vulnerabilities and corresponding fixes from open-source projects hosted on GitHub. CVEFixes, on the other hand, focuses on Common Vulnerabilities and Exposures (CVEs) and provides fine-grained fix commits associated with specific CVE entries. To maintain consistency with BigVul, we only retain C/C++ vulnerability samples from CVEFixes.We merge the two datasets following the standard practice established in prior work\cite{zhou2024large}, including deduplication and partitioning at the function level. Before filtering, the merged dataset consists of 135 CWEs, 4,067 CVEs, and 14,216 vulnerable functions. By applying the set of CWE identifiers identified in the previous step, we construct a refined dataset focused on overflow vulnerabilities, which includes 2,054 distinct CVEs and 7,865 vulnerable functions.A detailed breakdown of the final dataset is presented in Table~\ref{tab:cwe-distribution}.


\begin{table}[H]
\centering
\begin{tabular}{c|p{7.6cm}}
\hline
\textbf{ID} & \textbf{Name} \\
\hline
119 & Improper Restriction of Operations within the Bounds of a Memory Buffer \\
120 & Buffer Copy without Checking Size of Input ('Classic Buffer Overflow') \\
121 & Stack-based Buffer Overflow \\
122 & Heap-based Buffer Overflow \\
123 & Write-what-where Condition \\
124 & Buffer Underwrite (Buffer Underflow) \\
125 & Out-of-bounds Read \\
126 & Buffer Over-read \\
127 & Buffer Under-read \\
190$^\ast$ & Integer Overflow or Wraparound \\
415 & Double Free \\
416 & Use After Free \\
466 & Return of Pointer Value Outside of Expected Range \\
476$^\ast$ & NULL Pointer Dereference \\
761$^\ast$ & Free of Pointer not at Start of Buffer \\
786 & Access of Memory Location Before Start of Buffer \\
787 & Out-of-bounds Write \\
788 & Access of Memory Location After End of Buffer \\
805 & Buffer Access with Incorrect Length Value \\
806 & Buffer Access Using Size of Source Buffer \\
822 & Untrusted Pointer Dereference \\
823 & Use of Out-of-range Pointer Offset \\
824 & Access of Uninitialized Pointer \\
825 & Expired Pointer Dereference \\
\hline
\end{tabular}
\caption{List of CWE entries related to overflow vulnerabilities. Entries marked with $^\ast$ are added based on recent literature.}
\label{tab:cwe-list}
\end{table}

\begin{table}[H]
\centering

\begin{tabular}{c|c|c}
\hline
\textbf{CWE ID} & \textbf{CVE Count} & \textbf{Vulnerability Count} \\
\hline
CWE-119 & 483 & 1532 \\
CWE-120 & 71 & 249 \\
CWE-121 & 9 & 14 \\
CWE-122 & 47 & 69 \\
CWE-125 & 425 & 1417 \\
CWE-126 & 2 & 2 \\
CWE-190 & 159 & 2539 \\
CWE-415 & 42 & 162 \\
CWE-416 & 215 & 463 \\
CWE-476 & 267 & 456 \\
CWE-787 & 319 & 943 \\
CWE-823 & 4 & 5 \\
CWE-824 & 11 & 14 \\
\hline
\textbf{Total} & \textbf{2054} & \textbf{7865} \\
\hline
\end{tabular}
\caption{CWE Distribution}
\label{tab:cwe-distribution}
\end{table}

\subsection{RQ design}




\begin{itemize}
  \item \textbf{RQ1: Vulnerability Repairableness} 
  % Dataset: 614         7865 at confidence 99% and error 5%.
  %             无关的修复函数    单函数视角下无法修复   单函数视角下可修复 
                % Irrelevant Intra-Function Unrepairable   Intra-Function Repairable
  %单函数漏洞        N                  N                 Y
  %多函数漏洞        Y                  Y                 Y
  % 1. 分类的情况 2. Unrepairable的漏洞函数会如何影响"漏洞修复工具"的效果？  

  % Assume 首先修复Intra-Function Repairable之后，Intra-Function Unrepairable 的函数，增加上下文之后，是否可以变成可修复的？--------Y

  % 结论：有些漏洞函数的修复，是unrepairable的，不能作为评估repair的验证集

  \item \textbf{RQ2. Repairable Vulnerability Characterization.} What are the key characteristics of overflow vulnerabilities?
  % 可以repair的漏洞函数，他们的特征；不同类型和特征，
  \begin{itemize}
    \item \textbf{Trigger Logic.} What are the common triggering patterns or logic that lead to overflows?
    \item \textbf{Fixing Strategies.} What kinds of repair strategies are commonly applied to overflow vulnerabilities?
  \end{itemize}

  % Dataset: 600
  % 使用漏洞修复工具评估的效果和区别？
  \item \textbf{RQ3 LLM-based Repairable Vulnerability Repair Evaluation} Effectiveness with regard to baselines, vulnerability type and fixing strategies.
  
  % A. 统计层面的分析
  % baseline: 1. simple prompting 2. agent-based ... 找一下相关的工作
  % trigger logic. fixing strategies 
  % B. 人工采样的分析(定量分析)

  % 结论： 发现多agent比单agent效果好，过度修复的问题

\end{itemize}



% \begin{itemize}[leftmargin=*]
 
  
%   \item \textbf{RQ1. Effectiveness of Simple Prompting.} How effective is a simple prompt-based approach in enabling LLMs to repair overflow vulnerabilities?
% \begin{itemize}
%   \item \textbf{RQ1.1 Impact of Trigger Logic.} How does the effectiveness of prompt-based LLM repair vary across different overflow trigger logics?
%   \item \textbf{RQ1.2 Impact of Fixing Strategies.} How does the effectiveness of prompt-based LLM repair vary across different fixing strategies?
% \end{itemize}

% \item \textbf{RQ2. Effectiveness of Agent-based Repair.} How effective is an agent-assisted approach in enabling LLMs to repair overflow vulnerabilities?
% \begin{itemize}
%   \item \textbf{RQ2.1 Impact of Trigger Logic.} How does the effectiveness of agent-based LLM repair vary across different overflow trigger logics?
%   \item \textbf{RQ2.2 Impact of Fixing Strategies.} How does the effectiveness of agent-based LLM repair vary across different fixing strategies?
% \end{itemize}

% \end{itemize}

\subsection{RQ0: Vulnerability Characterization}

\subsubsection{ploit labeling}

xxxx

\begin{tcolorbox}[colback=gray!5!white, colframe=black!75, title=Answer to RQ0]

Overflow vulnerabilities can be characterized along two axes: the \textit{trigger logic} and the \textit{repair strategy}. 

\end{tcolorbox}


